\documentclass[12pt]{article}
\usepackage[utf8]{inputenc}
\usepackage[T1]{fontenc}
\usepackage[margin=1in]{geometry}
\usepackage{setspace}
\usepackage{xcolor}
\usepackage[colorlinks=true, urlcolor=blue!60!black, linkcolor=blue!60!black]{hyperref}
\setlength{\parindent}{0em}
\setlength{\parskip}{0.45em}
\renewcommand{\baselinestretch}{1.0}
\begin{document}
\begin{center}
{\Large \textit{Troxel v. Granville}, 530 U.S. 57 (2000)\par}
\vspace{0.3em}
{\normalsize Case Brief --- One-Page Edition\par}
\end{center}

\textbf{Memory Jogger.} A judge may not override a fit mother's limits on grandparent visitation under a ``breathtakingly broad'' best-interest statute (\href{https://www.law.cornell.edu/supremecourt/text/99-138}{Cornell LII}).

\textbf{Facts.} After their son died, Jenifer and Gary Troxel sought more time with his daughters than their mother, Tommie Granville, would allow; she offered one day a month. Washington let ``any person'' petition at any time for visitation serving a child's best interest. Though Granville was concededly fit, the trial court ordered a weekend a month and a week each summer, presuming visitation was best.

\textbf{Procedural History.} The Washington Supreme Court struck down the statute on its face for lacking a harm requirement.

\textbf{Votes.} Affirmed, 6--3, with no majority: O'Connor for four; Souter and Thomas concurred in the judgment; Stevens, Scalia, and Kennedy dissented.

\textbf{Holding.} Parents have a fundamental right to direct their children's care, custody, and control; courts must give a fit parent's decision special weight. The order was unconstitutional as applied; whether harm must be shown was left open.

\textbf{Reasoning.} Granville was fit, the court presumed against her, ignored her offer of visitation, and relied on thin findings---cousins, music, and the judge's own memories of his grandparents---showing only a disagreement about what was better for the girls. Souter would invalidate the statute facially; Thomas would apply strict scrutiny. Stevens stressed the child's own interests, Scalia refused to enforce unenumerated rights against state law, and Kennedy rejected a rigid harm rule.

\textbf{Examples.} \emph{Same-Side:} a law letting any neighbor obtain visits. \emph{Opposite-Side:} visits ordered after a parent is found unfit. \emph{Fence-Sitter:} a de facto parent's petition after a breakup.

\textbf{Critique.} The fractured opinions left parental rights ``adrift in the middle,'' with no standard of review and no definition of special weight; child-welfare scholars add that rigid parental autonomy slights children's bonds in blended and kinship families.

\textbf{Key Quotation.} ``The liberty interest at issue in this case---the interest of parents in the care, custody, and control of their children---is perhaps the oldest of the fundamental liberty interests recognized by this Court.'' --- Justice O'Connor, plurality opinion.

\end{document}
